
\chapter{Continuous morphisms}

The three forms of continuity, their equivalence for entirely defined
monovalued arrows, and their behavior under restriction are established
in Section~\ref{sec-osa-continuity}. Here I compare them with the
usual continuity conditions of particular spaces. For funcoids and
reloids the dagger of an arrow is its reverse.

\section{Traditional definitions of continuity}

In this section we will show that having a funcoid or reloid $\uparrow f$
corresponding to a function $f$ we can express continuity of it by
the formula $\uparrow f\circ\mu\sqsubseteq\nu\circ\uparrow f$ (or
similar formulas) where $\mu$ and $\nu$ are some spaces.


\subsection{\index{continuity!pre-topology}Pretopology}

Let $(A,\cl_{A})$ and $(B,\cl_{B})$ be preclosure spaces. Then by
definition a function $f:A\rightarrow B$ is continuous iff $f\cl_{A}(X)\subseteq\cl_{B}(fX)$
for every $X\in\subsets A$. Let now $\mu$ and $\nu$ be endofuncoids
corresponding correspondingly to $\cl_{A}$ and $\cl_{B}$. Then the
condition for continuity can be rewritten as
\[
\uparrow^{\mathsf{FCD}(\Ob\mu,\Ob\nu)}f\circ\mu\sqsubseteq\nu\circ\uparrow^{\mathsf{FCD}(\Ob\mu,\Ob\nu)}f.
\]



\subsection{\index{continuity!proximity}Proximity spaces}

Let $\mu$ and $\nu$ be proximity spaces (which I consider a special
case of endofuncoids). By definition a $\mathbf{Set}$-morphism~$f$
is a proximity-continuous map from $\mu$
to $\nu$ iff
\[
\forall X,Y\in\mathscr{T}(\Ob\mu):(X\rsuprel{\mu}Y\Rightarrow\rsupfun fX\rsuprel{\nu}\rsupfun fY).
\]


Equivalently transforming this formula we get
\begin{gather*}
\forall X,Y\in\mathscr{T}(\Ob\mu):(X\rsuprel{\mu}Y\Rightarrow\supfun f\uparrow X\suprel{\nu}\supfun f\uparrow Y);\\
\forall X,Y\in\mathscr{T}(\Ob\mu):(X\rsuprel{\mu}Y\Rightarrow\uparrow X\suprel{f^{-1}\circ\nu\circ f}\uparrow Y);\\
\forall X,Y\in\mathscr{T}(\Ob\mu):(X\rsuprel{\mu}Y\Rightarrow X\rsuprel{f^{-1}\circ\nu\circ f}Y);\\
\mu\sqsubseteq f^{-1}\circ\nu\circ f.
\end{gather*}


So a function $f$ is proximity continuous iff $\mu\sqsubseteq f^{-1}\circ\nu\circ f$.


\subsection{\index{continuity!uniformity}Uniform spaces}

Uniform spaces are a special case of endoreloids.

Let $\mu$ and $\nu$ be uniform spaces. By definition a $\mathbf{Set}$-morphism~$f$
is a uniformly continuous map from $\mu$ to $\nu$ iff 
\[
\forall\varepsilon\in\up\nu\exists\delta\in\up\mu\forall(x,y)\in\delta:(fx,fy)\in\varepsilon.
\]


Equivalently transforming this formula we get:
\begin{gather*}
\forall\epsilon\in\up\nu\exists\delta\in\up\mu\forall(x,y)\in\delta:\{(fx,fy)\}\subseteq\epsilon;\\
\forall\epsilon\in\up\nu\exists\delta\in\up\mu\forall(x,y)\in\delta:f\circ\{(x,y)\}\circ f^{-1}\subseteq\epsilon;\\
\forall\epsilon\in\up\nu\exists\delta\in\up\mu:f\circ\delta\circ f^{-1}\subseteq\epsilon;\\
\forall\epsilon\in\up\nu:\uparrow^{\mathsf{RLD}(\Ob\mu,\Ob\nu)}f\circ\mu\circ(\uparrow^{\mathsf{RLD}(\Ob\mu,\Ob\nu)}f)^{-1}\sqsubseteq\uparrow^{\mathsf{RLD}(\Ob\mu,\Ob\nu)}\epsilon;\\
\uparrow^{\mathsf{RLD}(\Ob\mu,\Ob\nu)}f\circ\mu\circ(\uparrow^{\mathsf{RLD}(\Ob\mu,\Ob\nu)}f)^{-1}\sqsubseteq\nu.
\end{gather*}
So a function $f$ is uniformly continuous iff $f\circ\mu\circ f^{-1}\sqsubseteq\nu$.


\section{Passage from reloids to funcoids}

The classical general topology theorem that uniformly continuous function
from a uniform space to an other uniform space is proximity-continuous
regarding the proximities generated by the uniformities, generalized
for reloids and funcoids takes the following form:
\begin{thm}
If an entirely defined morphism of the category of reloids $f\in\continuous''(\mu,\nu)$
for some endomorphisms $\mu$ and $\nu$ of the category of reloids,
then $\mathsf{\tofcd}f\in\continuous'(\tofcd\mu,\tofcd\nu)$.\end{thm}
\begin{xca}
I leave a simple exercise for the reader to prove the last theorem.
\end{xca}

\section{Continuity for topological spaces}

\begin{prop}
  The following are pairwise equivalent for funcoids $\mu$, $\nu$ and
  a monovalued, entirely defined morphism~$f\in\Hom(\Ob\mu,\Ob\nu)$:
  \begin{enumerate}
    \item\label{fcd-cont-unfold} $\forall A \in \mathscr{T} \Ob \mu, B \in \up
    \supfun{\nu} \rsupfun{f} A : \rsupfun{f^{- 1}}
    B \in \up \rsupfun{\mu} A$.

    \item\label{fcd-cont-main}  $f \in \continuous (\mu, \nu)$.

    \item\label{fcd-cont-rev} $f \in \continuous (\mu^{- 1} , \nu^{- 1})$.
  \end{enumerate}
\end{prop}

\begin{proof}
  ~
  \begin{description}
  \item[\ref{fcd-cont-main}$\Leftrightarrow$\ref{fcd-cont-rev}] By general $f \circ \mu \sqsubseteq \nu \circ
  f \Leftrightarrow f \circ \mu^{\dagger} \sqsubseteq \nu^{\dagger}
  \circ f$ formula above.

  \item[\ref{fcd-cont-unfold}$\Leftrightarrow$\ref{fcd-cont-main}] \ref{fcd-cont-unfold} is equivalent to
  $\rsupfun{\rsupfun{f^{- 1}}}
  \up \supfun{\nu} \rsupfun{f} A \subseteq \up
  \rsupfun{\mu} A$ equivalent to $\supfun{\nu} \rsupfun{f}
  A \sqsupseteq \supfun{f} \rsupfun{\mu} A$
  (used ``Orderings of filters'' chapter).
  \end{description}
\end{proof}

\begin{cor}\label{top-cont}
  The following are pairwise equivalent for topological spaces $\mu$,
  $\nu$ and a monovalued, entirely defined morphism~$f\in\Hom(\Ob\mu,\Ob\nu)$:
  \begin{enumerate}
    \item\label{top-cont-unfold} $\forall x \in \Ob \mu, B \in \up \supfun{\nu}
    \rsupfun{f} \{x\} : \rsupfun{f^{- 1}} B \in
    \up \rsupfun{\mu} \{x\}$.

    \item\label{top-cont-preim} Preimages (by $f$) of open sets are open.

    \item\label{top-cont-main} $f \in \continuous (\mu, \nu)$ that is $\supfun{f}
    \rsupfun{\mu} \{ x \} \sqsubseteq \supfun{\nu}
    \rsupfun{f} \{ x \}$ for every $x \in \Ob \mu$.

    \item\label{top-cont-rev} $f \in \continuous (\mu^{- 1} , \nu^{- 1})$ that is $\supfun{f}
    \rsupfun{\mu^{- 1}} A \sqsubseteq \supfun{\nu^{- 1}}
    \rsupfun{f} A$ for every $A \in \mathscr{T}
    \Ob \mu$.
  \end{enumerate}
\end{cor}

\begin{proof}
  \ref{fcd-cont-main} from the previous proposition is equivalent to
  $\supfun{f}\rsupfun{\mu}\{x\}\sqsubseteq\supfun{\nu}\rsupfun{f}\{x\}$
  equivalent to $\rsupfun{\rsupfun{f^{- 1}}}
  \up \supfun{\nu} \rsupfun{f} \{ x \} \subseteq
  \up \rsupfun{\mu} \{ x \}$ for every $x \in \Ob
  \mu$, equivalent to \ref{top-cont-unfold} (used ``Orderings of filters'' chapter).

  It remains to prove \ref{top-cont-main}$\Leftrightarrow$\ref{top-cont-preim}.

  \begin{description}
  \item[\ref{top-cont-main}$\Rightarrow$\ref{top-cont-preim}] Let $B$ be an open set in $\nu$. For every $x \in
  \rsupfun{f^{- 1}} B$ we have $f (x) \in B$ that is $B$ is a
  neighborhood of $f (x)$, thus $\rsupfun{f^{-1}}B$ is a neighborhood of $x$. We have
  proved that $\rsupfun{f^{- 1}} B$ is open.

  \item[\ref{top-cont-preim}$\Rightarrow$\ref{top-cont-main}] Let $B$ be a neighborhood of $f (x)$. Then there is an
  open neighborhood $B' \subseteq B$ of $f (x)$. $\rsupfun{f^{-1}}
  B'$ is open and thus is a neighborhood of $x$ ($x \in \rsupfun{f^{-1}} B'$ because $f (x) \in B'$). Consequently
  $\rsupfun{f^{-1}} B$ is a neighborhood of $x$.
  \end{description}

  Alternative proof of \ref{top-cont-preim}$\Leftrightarrow$\ref{top-cont-rev}:
  \url{http://math.stackexchange.com/a/1855782/4876}
\end{proof}

\section{\texorpdfstring{$\continuous (\mu \circ \mu^{- 1} , \nu \circ \nu^{- 1})$}{C(mu o mu\textasciicircum-1, nu o nu\textasciicircum-1)}}

\begin{prop}
  $f \in \continuous (\mu, \nu) \Rightarrow f \in \continuous'' (\mu
  \circ \mu^{- 1} , \nu \circ \nu^{- 1})$ for endofuncoids $\mu$,
  $\nu$ and monovalued funcoid $f \in \mathsf{FCD} (\Ob
  \mu, \Ob \nu)$.
\end{prop}

\begin{proof}
  Let $f \in \continuous (\mu, \nu)$.

\begin{multline*}
  X \mathrel{[f \circ \mu \circ \mu^{- 1} \circ f^{- 1}]^{\ast}}
  Z \Leftrightarrow \\ \exists p \in \atoms^{\mathscr{F}} : \left( X
  \mathrel{[\mu^{- 1} \circ f^{- 1}]^{\ast}} p \wedge p \mathrel{[f
  \circ \mu]^{\ast}} Z \right) \Leftrightarrow \\ \exists p \in
  \atoms^{\mathscr{F}} : \left( p \mathrel{[f \circ \mu]^{\ast}} X
  \wedge p \mathrel{[f \circ \mu]^{\ast}} Z \right) \Rightarrow \\ \exists
  p \in \atoms^{\mathscr{F}} : \left( p \mathrel{[\nu \circ f]^{\ast}} X
  \wedge p \mathrel{[\nu \circ f]^{\ast}} Z \right) \Leftrightarrow \\ \exists p
  \in \atoms^{\mathscr{F}} : \left( \rsupfun{f} p
  \mathrel{[\nu]^{\ast}} X \wedge \rsupfun{f} p
  \mathrel{[\nu]^{\ast}} Z \right) \Rightarrow X \mathrel{[\nu \circ \nu^{-
  1}]^{\ast}} Z
\end{multline*}
(taken into account monovaluedness of $f$ and thus that
  $\rsupfun{f} p$ is atomic or least). Thus $f \circ \mu
  \circ \mu^{- 1} \circ f^{- 1} \sqsubseteq \nu \circ \nu^{- 1}$ that is
  $f \in \continuous'' (\mu \circ \mu^{- 1} , \nu \circ \nu^{-
  1})$.
\end{proof}

\begin{prop}
  $f \in \continuous'' (\mu \circ \mu^{- 1} , \nu \circ \nu^{- 1})
  \Rightarrow f \in \continuous'' (\mu, \nu)$ for complete endofuncoids
  $\mu$, $\nu$ and principal funcoid $f \in \mathsf{FCD}
  (\Ob \mu, \Ob \nu)$, provided that $\mu$ is
  reflexive, and $\nu$ is $T_1$-separable.
\end{prop}

\begin{proof}
\begin{multline*}
  f \in \continuous'' (\mu \circ \mu^{- 1} , \nu \circ \nu^{- 1})
  \Leftrightarrow \\ f \circ \mu \circ \mu^{- 1} \circ f^{- 1}
  \sqsubseteq \nu \circ \nu^{- 1} \Rightarrow \text{(reflexivity of
  $\mu$)} \Rightarrow \\ f \circ \mu \circ f^{- 1} \sqsubseteq \nu
  \circ \nu^{- 1} \Leftrightarrow f \circ \mu^{- 1} \circ f^{- 1}
  \sqsubseteq \nu \circ \nu^{- 1} \Rightarrow \\ \langle f \circ \mu^{- 1}
  \circ f^{- 1} \rangle^{\ast} X \sqsubseteq \supfun{\nu}^{\ast}
  \langle \nu^{- 1} \rangle^{\ast} X \Rightarrow \\ \Cor \left\langle f
  \circ \mu^{- 1} \circ f^{- 1} \right\rangle^{\ast} X \sqsubseteq
  \Cor \supfun{\nu}^{\ast} \langle \nu^{- 1} \rangle^{\ast} X
  \Leftrightarrow \\ \langle f \circ \mu^{- 1} \circ f^{- 1} \rangle^{\ast}
  X \sqsubseteq \Cor \supfun{\nu}^{\ast} \langle \nu^{- 1}
  \rangle^{\ast} X \Rightarrow \\ \text{($T_1$-separability)} \Rightarrow \\ \langle
  f \circ \mu^{- 1} \circ f^{- 1} \rangle^{\ast} X \sqsubseteq \langle
  \nu^{- 1} \rangle^{\ast} X$ for any typed set $X$ on $\Ob \nu.
\end{multline*}
  Thus
  \begin{multline*}
f \in \continuous'' (\mu \circ \mu^{- 1} , \nu \circ \nu^{- 1})
  \Rightarrow f \circ \mu^{- 1} \circ f^{- 1} \sqsubseteq \nu^{- 1}
  \Leftrightarrow \\ f \circ \mu \circ f^{- 1} \sqsubseteq \nu
  \Leftrightarrow f \in \continuous'' (\mu, \nu).
  \end{multline*}
\end{proof}

\begin{thm}
  $f \in \continuous (\mu \circ \mu^{- 1} , \nu \circ \nu^{- 1})
  \Leftrightarrow f \in \continuous (\mu, \nu)$ for complete endofuncoids
  $\mu$, $\nu$ and principal monovalued and entirely defined funcoid $f
  \in \mathsf{FCD} (\Ob \mu, \Ob \nu)$, provided that
  $\mu$ is reflexive, and $\nu$ is $T_1$-separable.
\end{thm}

\begin{proof}
  Two above propositions and theorem~\ref{cont-eq}.
\end{proof}

\section{Anticontinuity and open maps}

\begin{defn}
\index{open map}An \emph{open map} from a topological space to a
topological space is a function which maps open sets into open sets.
\end{defn}

\begin{thm}
For topological spaces considered as complete funcoids,
a principal anticontinuous morphism is the same as open map.
\end{thm}

\begin{proof}
Because $f$,~$\mu$,~$\nu$ are complete funcoids, we have
\[
f\in\continuous_{\ast}(\mu,\nu)\Leftrightarrow
f\circ\mu\sqsupseteq\nu\circ f\Leftrightarrow
\Compl(f\circ\mu)\sqsupseteq\Compl(\nu\circ f).
\]
Equivalently transforming further, we get
\[
\forall x\in\Ob\mu:\supfun f\rsupfun{\mu}@\{x\}\sqsupseteq\supfun{\nu}\rsupfun
f@\{x\};
\]
\[
\forall x\in\Ob\mu,V\in\rsupfun{\mu}\{x\}:\rsupfun
fV\sqsupseteq\supfun{\nu}\rsupfun f@\{x\},
\]
what is the criterion of~$f$ being an open map.
\end{proof}